%==============================================================
%  lecons/analyse/interversion_somme_integrale.tex
%  Théorèmes d'interversion somme — intégrale
%==============================================================

\lecontitle{Interversion somme et intégrale}{Analyse réelle — Séries de fonctions et intégration}

%=============================================================
%  § 1 — ÉNONCÉS
%=============================================================
\secnum{navyblue}{1}{Énoncés}

\begin{tcolorbox}[thmbox]
  \textbf{Problème.}\enspace%
  \index{interversion somme--intégrale}
  Étant donnée une série de fonctions $\sum_{n\geq 0} f_n$ convergeant
  (simplement) vers $f$ sur un intervalle $I$, sous quelles hypothèses
  peut-on écrire
  \[
    \int_I \sum_{n=0}^{+\infty} f_n(t)\,\mathrm{d}t
    \;=\; \sum_{n=0}^{+\infty} \int_I f_n(t)\,\mathrm{d}t\;\;?
  \]
  La question est non triviale dès que l'intervalle est non borné ou que
  les $f_n$ ne sont pas uniformément contrôlées.

  \medskip
  \textbf{Théorème 1 (convergence uniforme sur un segment).}\enspace%
  \index{convergence uniforme}\index{convergence normale}
  Soit $(f_n)$ une suite de fonctions continues sur un segment $[a,b]$
  convergeant \emph{uniformément} vers $f$. Alors $f$ est continue et
  \[
    \int_a^b f_n(t)\,\mathrm{d}t \;\xrightarrow[n\to+\infty]{}\;
    \int_a^b f(t)\,\mathrm{d}t.
  \]
  En particulier, si $\sum f_n$ converge \emph{normalement} sur $[a,b]$,
  \[
    \int_a^b \sum_{n=0}^{+\infty} f_n(t)\,\mathrm{d}t
    \;=\; \sum_{n=0}^{+\infty} \int_a^b f_n(t)\,\mathrm{d}t.
  \]

  \medskip
  \textbf{Théorème 2 (sommation des séries à termes positifs).}\enspace%
  \index{Beppo Levi (théorème de)}
  Soit $(f_n)$ une suite de fonctions \emph{positives} et continues par
  morceaux sur un intervalle $I$. Alors, dans $[0,+\infty]$~:
  \[
    \int_I \sum_{n=0}^{+\infty} f_n(t)\,\mathrm{d}t
    \;=\; \sum_{n=0}^{+\infty} \int_I f_n(t)\,\mathrm{d}t.
  \]
  La somme $\sum f_n$ est intégrable sur $I$ si et seulement si
  $\sum_n \int_I f_n < +\infty$.

  \medskip
  \textbf{Théorème 3 (intégration terme à terme).}\enspace%
  \index{intégration terme à terme}
  Soit $(f_n)$ une suite de fonctions continues par morceaux et
  intégrables sur $I$. On suppose~:
  \begin{enumerate}
    \item la série $\sum f_n$ converge simplement vers une fonction $f$
          continue par morceaux sur $I$~;
    \item $\displaystyle\sum_{n=0}^{+\infty}\int_I |f_n(t)|\,\mathrm{d}t
          \;<\; +\infty.$
  \end{enumerate}
  Alors $f$ est intégrable sur $I$ et
  \[
    \int_I f(t)\,\mathrm{d}t
    \;=\; \sum_{n=0}^{+\infty} \int_I f_n(t)\,\mathrm{d}t.
  \]
\end{tcolorbox}

\begin{significationintuitive}
Permuter $\sum$ et $\int$ revient à échanger deux limites~: l'intégrale est
un passage à la limite (sommes de Riemann), et la série en est un autre
(sommes partielles). Une telle interversion exige toujours une forme
d'\emph{uniformité} ou de \emph{domination}. Trois cadres se dégagent~:
sur un segment, la convergence uniforme suffit~; sur un intervalle
quelconque, la positivité (Théorème~2) ou la sommabilité absolue
(Théorème~3) prennent le relais.
\end{significationintuitive}

\decsep{}

%=============================================================
%  § 2 — DÉMONSTRATION
%=============================================================
\secnum{navyblue}{2}{Démonstrations}

\rubriq{Théorème 1 — convergence uniforme}

\begin{mdframed}[style=proofbar]
  \textit{Continuité de la limite.}\enspace%
  La limite uniforme d'une suite de fonctions continues est continue
  (résultat classique)~; en particulier $f$ est intégrable sur $[a,b]$.

  \medskip
  \textit{Passage à la limite.}\enspace%
  La majoration fondamentale de l'intégrale sur un segment donne
  \[
    \left|\int_a^b f_n - \int_a^b f\right|
    \;\leq\; \int_a^b |f_n(t) - f(t)|\,\mathrm{d}t
    \;\leq\; (b-a)\,\|f_n - f\|_{\infty,[a,b]}.
  \]
  Comme $\|f_n - f\|_\infty \to 0$ par convergence uniforme, le membre
  de gauche tend vers $0$.

  \medskip
  \textit{Application aux séries normalement convergentes.}\enspace%
  Soit $S_N = \sum_{n=0}^{N} f_n$. La convergence normale entraîne la
  convergence uniforme de $(S_N)$ vers $S = \sum f_n$ sur $[a,b]$. On
  applique le résultat précédent à $(S_N)$~:
  \[
    \int_a^b S_N \;=\; \sum_{n=0}^{N}\int_a^b f_n
    \;\xrightarrow[N\to+\infty]{}\;
    \int_a^b S \;=\; \int_a^b \sum_{n=0}^{+\infty} f_n.
  \]
  \hfill$\blacksquare$
\end{mdframed}

\rubriq{Théorème 2 — termes positifs}

\begin{mdframed}[style=proofbar]
  Notons $S_N = \sum_{n=0}^{N} f_n$ et $S = \sum_{n=0}^{+\infty} f_n$
  (avec valeur éventuellement infinie). Les $f_n$ étant positives, la
  suite $(S_N)$ est croissante et tend simplement vers $S$ sur $I$.

  \medskip
  \textit{Inégalité $(\leq)$.}\enspace%
  Pour tout $N$, $S_N \leq S$, donc par croissance de l'intégrale
  des fonctions positives,
  \[
    \sum_{n=0}^{N}\int_I f_n \;=\; \int_I S_N \;\leq\; \int_I S.
  \]
  En passant à la limite $N\to+\infty$~: $\sum_{n=0}^{+\infty}\int_I f_n \leq \int_I S$.

  \medskip
  \textit{Inégalité $(\geq)$.}\enspace%
  Soit $[a,b]\subset I$ un segment quelconque. Par positivité,
  $\int_a^b S_N \leq \sum_{n=0}^{N}\int_I f_n$. La suite $(S_N)$ est
  croissante de limite $S$ sur le segment $[a,b]$, et le passage
  $\int_a^b S_N \to \int_a^b S$ relève du théorème de convergence
  monotone (Beppo Levi) appliqué sur ce segment. Donc
  \[
    \int_a^b S \;\leq\; \sum_{n=0}^{+\infty}\int_I f_n.
  \]
  En prenant l'enveloppe supérieure sur les segments $[a,b]\subset I$,
  on obtient $\int_I S \leq \sum_{n=0}^{+\infty}\int_I f_n$.

  \medskip
  Les deux inégalités donnent l'égalité (avec valeur $+\infty$
  autorisée des deux côtés).\hfill$\blacksquare$
\end{mdframed}

\rubriq{Théorème 3 — intégration terme à terme}

\begin{mdframed}[style=proofbar]
  \textit{Intégrabilité de $f$.}\enspace%
  La fonction $|f| = \bigl|\sum f_n\bigr| \leq \sum |f_n|$ est dominée
  par une fonction positive dont l'intégrale vaut, par le Théorème~2,
  $\sum_n \int_I |f_n| < +\infty$. Donc $f$ est intégrable sur $I$.

  \medskip
  \textit{Convergence de $\int_I (f - S_N)$ vers $0$.}\enspace%
  Le reste $f - S_N = \sum_{k>N} f_k$ vérifie
  \[
    |f - S_N| \;\leq\; \sum_{k=N+1}^{+\infty} |f_k|.
  \]
  Le Théorème~2 appliqué aux $|f_k|$ donne
  \[
    \int_I |f - S_N| \;\leq\;
    \sum_{k=N+1}^{+\infty} \int_I |f_k|
    \;\xrightarrow[N\to+\infty]{}\; 0,
  \]
  comme reste d'une série convergente à termes positifs.

  \medskip
  \textit{Conclusion.}\enspace%
  Par l'inégalité triangulaire,
  \[
    \left|\int_I f \;-\; \sum_{n=0}^{N}\int_I f_n\right|
    \;=\; \left|\int_I (f - S_N)\right|
    \;\leq\; \int_I |f - S_N|
    \;\xrightarrow[N\to+\infty]{}\; 0.
  \]
  Ainsi $\sum_n \int_I f_n = \int_I f$.\hfill$\blacksquare$
\end{mdframed}

\decsep{}

%=============================================================
%  § 3 — EXEMPLES, CONTRE-EXEMPLES, EXERCICES
%=============================================================
\secnum{exgreen}{3}{Exemples, contre-exemples et exercices}

\rubriq{Exemples}

\begin{mdframed}[style=exbar]
  \textbf{1. Développement en série de $\ln(1+x)$.}\enspace%
  Pour $x \in {]-1,1[}$, $\dfrac{1}{1+t} = \sum_{n=0}^{+\infty} {(-1)}^n t^n$
  uniformément sur tout segment $[-r,r]$ avec $r<1$. Une intégration terme à terme
  donne
  \[
    \ln(1+x) \;=\; \sum_{n=0}^{+\infty}\frac{{(-1)}^n}{n+1}x^{n+1}
    \;=\; \sum_{n=1}^{+\infty}\frac{{(-1)}^{n-1}}{n}x^n.
  \]

  \smallskip\noindent
  \textbf{2. Fonction zêta et intégrale.}\enspace%
  Pour $s > 1$,
  $\displaystyle\zeta(s)\,\Gamma(s) = \int_0^{+\infty}\frac{t^{s-1}}{e^t-1}\,\mathrm{d}t$.
  La preuve consiste à écrire $\dfrac{1}{e^t - 1} = \sum_{n\geq 1} e^{-nt}$
  (série à termes positifs), puis à appliquer le Théorème~2~:
  $\int_0^{+\infty} t^{s-1}\sum_{n\geq 1} e^{-nt}\,\mathrm{d}t
  = \sum_{n\geq 1}\int_0^{+\infty} t^{s-1}e^{-nt}\,\mathrm{d}t
  = \sum_{n\geq 1}\Gamma(s)/n^s$.

  \smallskip\noindent
  \textbf{3. Identité d'Euler $\sum 1/n^2 = \pi^2/6$.}\enspace%
  Le développement
  $\dfrac{1}{1-tu} = \sum_{n\geq 0}{(tu)}^n$ et le Théorème~2 donnent
  $\displaystyle\int_0^1\!\int_0^1 \frac{\mathrm{d}t\,\mathrm{d}u}{1-tu}
   = \sum_{n=0}^{+\infty}\frac{1}{{(n+1)}^2} = \zeta(2)$.
  Un changement de variable astucieux permet alors de calculer
  l'intégrale et de retrouver $\pi^2/6$ (Beukers, Calabi).

  \smallskip\noindent
  \textbf{4. Formule des compléments.}\enspace%
  $\displaystyle\int_0^{+\infty} \frac{t^{s-1}}{1+t}\,\mathrm{d}t
  = \frac{\pi}{\sin(\pi s)}$ pour $0<s<1$. Le développement
  $\dfrac{1}{1+t} = \sum {(-1)}^n t^n$ ne vaut que sur $[0,1[$~; on découpe
  donc l'intégrale en $\int_0^1 + \int_1^{+\infty}$, on intègre terme à
  terme sur $[0,1]$ (Théorème~3), puis le changement $t \mapsto 1/t$
  ramène la partie $[1,+\infty[$ à une intégrale sur $[0,1]$ traitée de même.
\end{mdframed}

\vspace{6pt}
\rubriq{Contre-exemples et pièges}

\begin{mdframed}[style=cexbar]
  \textbf{Sans uniformité, l'interversion peut échouer.}\enspace%
  Soit $f_n : [0,1] \to \mathbb{R}$ définie par $f_n(t) = n$ si
  $t\in{]0,1/n]}$, et $0$ sinon. On a $\int_0^1 f_n = 1$ pour tout $n$,
  alors que $f_n(t) \to 0$ pour tout $t \in [0,1]$. Donc
  $\int \lim f_n = 0 \neq 1 = \lim \int f_n$.
  En posant $g_n = f_n - f_{n-1}$ (avec $f_{-1}=0$), la série $\sum g_n$
  donne un contre-exemple à l'interversion.

  \smallskip\noindent
  \textbf{L'hypothèse de positivité est essentielle pour le Théorème~2.}\enspace%
  Sans positivité, des compensations entre termes peuvent rendre les
  deux membres différents (ou rendre l'un convergent et l'autre pas).
  L'exemple précédent, lu comme série télescopique signée, illustre
  ce phénomène.

  \smallskip\noindent
  \textbf{$\sum\int |f_n| < +\infty$ est strictement plus forte que la
  convergence simple.}\enspace%
  La série $\sum_{n\geq 1}\dfrac{{(-1)}^n}{n}\mathbf{1}_{[n,n+1]}$
  converge simplement sur $[1,+\infty[$ vers une fonction dont la somme est
  intégrable (sur chaque palier $[n,n+1]$ elle vaut $(-1)^n/n$, et la série
  $\sum (-1)^n/n$ converge), mais $\sum\int |f_n| = \sum 1/n = +\infty$~;
  l'interversion via le Théorème~3 ne s'applique pas (bien que, par
  d'autres arguments, le résultat puisse rester valable).

  \smallskip\noindent
  \textbf{Convergence simple $\neq$ convergence uniforme.}\enspace%
  Sur $[0,1]$, $f_n(t) = t^n$ converge simplement vers $\mathbf{1}_{\{1\}}$
  mais non uniformément. Pourtant l'interversion $\int f_n = 1/(n+1)\to 0$
  fonctionne ici~; la convergence uniforme est suffisante, pas nécessaire.
\end{mdframed}

\vspace{6pt}
\exolabel

\begin{mdframed}[style=exobar]
  \begin{enumerate}
    \item Montrer que $\displaystyle\int_0^1\frac{\ln t}{t-1}\,\mathrm{d}t
          = \sum_{n=1}^{+\infty}\frac{1}{n^2} = \frac{\pi^2}{6}$.
          (\textit{Indication}~: développer $1/(1-t)$ en série
          géométrique, puis appliquer le Théorème~2 après vérification
          du signe.)

    \item Calculer $\displaystyle\int_0^{+\infty}\frac{t}{e^t-1}\,\mathrm{d}t$
          et plus généralement, pour $s>1$,
          $\displaystyle\int_0^{+\infty}\frac{t^{s-1}}{e^t-1}\,\mathrm{d}t
          = \Gamma(s)\,\zeta(s)$.

    \item Montrer que $\displaystyle\int_0^1 t^{-t}\,\mathrm{d}t
          = \sum_{n=1}^{+\infty}\frac{1}{n^n}$
          (\emph{sophomore's dream}). \textit{Indication}~: écrire
          $t^{-t} = e^{-t\ln t} = \sum_n \frac{{(-t\ln t)}^n}{n!}$ puis
          intégrer terme à terme~; calculer $\int_0^1 t^n(-\ln t)^n\,\mathrm{d}t$
          par le changement $t = e^{-u}$.

    \item Soit $f$ continue sur $[0,1]$. Montrer que
          $\displaystyle\sum_{n=1}^{+\infty}\int_0^1 t^n f(t)\,\mathrm{d}t$
          converge et calculer sa somme en fonction de $f$.

    \item (\textit{Échec contrôlé}) Soit $f_n(t) = n\,t^n(1-t)$ sur $[0,1]$.
          Montrer que $f_n \to 0$ simplement mais non uniformément~;
          calculer $\int_0^1 f_n$ et conclure si l'interversion
          $\lim\int = \int\lim$ tient.

    \item En écrivant $\dfrac{1}{1-x} = \sum x^n$, retrouver
          $\displaystyle\int_0^1 \frac{-\ln(1-x)}{x}\,\mathrm{d}x = \zeta(2)$.

    \item Démontrer la formule d'Abel--Plana sous une forme simplifiée~:
          si $f$ est analytique et décroît assez vite,
          $\displaystyle\sum_{n=0}^{+\infty} f(n) = \int_0^{+\infty} f(t)\,\mathrm{d}t
          + \frac{f(0)}{2} + \text{terme correctif}$. Quelle interversion
          est en jeu~?
  \end{enumerate}
\end{mdframed}

\leconfooter{B.\ Levi (1906) — H.\ Lebesgue (1904) — G.\ Fubini (1907)}
